\RequirePackage{fix-cm}
\documentclass[12pt]{extarticle}
\newcommand{\bandlabel}{Grades 4--5}
% Shared preamble for the Week 1 student pages. \bandlabel is defined by each packet.
\usepackage[letterpaper,top=0.45in,bottom=0.45in,left=0.5in,right=0.5in,
  includehead,includefoot,headheight=18pt,headsep=12pt,footskip=24pt]{geometry}
\usepackage[T1]{fontenc}

\usepackage[default]{sourcesanspro}
\usepackage{tikz}
\usepackage{fancyhdr}
\usepackage{array}

\definecolor{pbgreen}{RGB}{110,195,110}
\definecolor{pbblue}{RGB}{95,150,225}
\definecolor{pbred}{RGB}{232,95,85}
\definecolor{pbyellow}{RGB}{250,212,55}
\definecolor{pbpurple}{RGB}{160,115,205}

\pagestyle{fancy}
\fancyhf{}
\lhead{Week 1 / Tiling / \bandlabel}
\lfoot{Bellingham Math Circle / Week 1}
\rfoot{\thepage}
\renewcommand{\headrulewidth}{0pt}
\renewcommand{\footrulewidth}{0pt}

\setlength{\parindent}{0pt}
\setlength{\parskip}{0pt}
\raggedright

\newcounter{prob}
\newcommand{\prob}{\stepcounter{prob}\textbf{Problem \theprob:}\ }
\newcommand{\fig}[1]{\input{figs/#1.tex}\unskip}
\newcommand{\icon}[1]{$\vcenter{\hbox{\input{figs/icon_#1.tex}}}$}
\newcommand{\bigicon}[1]{\raisebox{-0.15em}{\input{figs/iconbig_#1.tex}}}

% board with a letter label underneath
\newcommand{\board}[2]{\begin{tabular}[b]{@{}c@{}}\fig{#2}\\[3pt]\textbf{#1}\end{tabular}}

% ruled writing lines: \writelines{number}{width}
\newcommand{\writelines}[2]{\begin{tikzpicture}
  \foreach \i in {1,...,#1}{\draw[black!35,line width=0.5pt] (0.5pt,-\i*0.42in) -- ([xshift=-1pt]#2,-\i*0.42in);}
  \path (0,0) -- (0,-#1*0.42in-0.05in);
\end{tikzpicture}}

% answer box
\newcommand{\abox}{\tikz[baseline=0.18in]\draw[line width=0.9pt] (0,0) rectangle (0.8in,0.5in);}

% figure with an answer box underneath
\newcommand{\figbox}[1]{\begin{tabular}[b]{@{}c@{}}\fig{#1}\\[8pt]\abox\end{tabular}}

% arrows between pictures
\newcommand{\twoarrow}{\tikz[baseline=-0.5ex]\draw[<->,>=stealth,line width=1.3pt] (0,0)--(0.6in,0);}
\newcommand{\onearrow}{\tikz[baseline=-0.5ex]\draw[->,>=stealth,line width=1.2pt] (0,0)--(0.28in,0);}

\AtBeginDocument{\fontsize{13pt}{16.5pt}\selectfont}

\begin{document}

Use only the small blocks. On every board in this packet, each block must cover whole small triangles of the board. Blocks may not overlap or stick out past the thick outline.

\vspace{10pt}
\prob Cover the board below with blue rhombuses \icon{B}. Find all the different ways to do it, and draw each way on one of the small boards. Explain how you know you have found them all.

\vspace{14pt}
\begin{center}
\fig{g45_b122}\hspace{0.35in}%
\begin{tabular}[b]{@{}c@{\hspace{0.2in}}c@{\hspace{0.2in}}c@{}}
\fig{g45_b122_small} & \fig{g45_b122_small} & \fig{g45_b122_small}\\[12pt]
\fig{g45_b122_small} & \fig{g45_b122_small} & \fig{g45_b122_small}
\end{tabular}

\vspace{14pt}
\fig{g45_b122_small}\hspace{0.45in}\fig{g45_b122_small}\hspace{0.45in}\fig{g45_b122_small}\hspace{0.45in}\fig{g45_b122_small}
\end{center}

\vspace{0pt}
\writelines{4}{\linewidth}

\newpage
\prob When three blue rhombuses in a covering make a hexagon the size of a yellow~block~\icon{Y}, you can lift them off and put them back the other way, as in the picture.

\vspace{10pt}
\begin{center}
\fig{g45_flipA}\hspace{0.3in}\raisebox{0.48in}{\twoarrow}\hspace{0.3in}\fig{g45_flipB}
\end{center}

\vspace{6pt}
Use this move to change one covering of the board from Problem 1 into another. Which two coverings need the most moves to get from one to the other? How many moves do they need?

\vfill
\writelines{3}{\linewidth}

\newpage
\prob The move in Problem 2 is called a flip. Put blue rhombuses on the hexagon board to make covering A. Change it into covering B using only flips, with as few flips as you can. Do the same for C and D, and for E and F. Write the smallest number of flips you found in each box.

\vspace{16pt}
\begin{minipage}[t]{4.05in}
\vspace{0pt}
\fig{g45_hex2}
\end{minipage}%
\begin{minipage}[t]{3.45in}
\vspace{0pt}
\centering
\board{A}{g45_pair_A}\hspace{0.07in}\raisebox{0.75in}{\onearrow}\hspace{0.07in}\board{B}{g45_pair_B}\\[8pt]
\abox\\[22pt]
\board{C}{g45_pair_C}\hspace{0.07in}\raisebox{0.75in}{\onearrow}\hspace{0.07in}\board{D}{g45_pair_D}\\[8pt]
\abox\\[22pt]
\board{E}{g45_pair_E}\hspace{0.07in}\raisebox{0.75in}{\onearrow}\hspace{0.07in}\board{F}{g45_pair_F}\\[8pt]
\abox
\end{minipage}

\newpage
\prob Make covering A from Problem 3 on the hexagon board. Make flips, one at a time, until you are back at covering A. Find a way to do this with 4 flips that never flips the same hexagon twice in a row. Draw the covering after each flip on the small hexagons. Can you get back to covering A with 5 flips? With 7 flips?

\vspace{16pt}
\begin{center}
\fig{g45_hex2_tinyA}\hspace{0.05in}\raisebox{0.45in}{\onearrow}\hspace{0.05in}%
\fig{g45_hex2_tiny}\hspace{0.05in}\raisebox{0.45in}{\onearrow}\hspace{0.05in}%
\fig{g45_hex2_tiny}\hspace{0.05in}\raisebox{0.45in}{\onearrow}\hspace{0.05in}%
\fig{g45_hex2_tiny}\hspace{0.05in}\raisebox{0.45in}{\onearrow}\hspace{0.05in}%
\fig{g45_hex2_tiny}
\end{center}

\vspace{10pt}
\writelines{6}{\linewidth}

\newpage
\prob In a covered board, start at a short edge on the top side. Go down through the rhombuses, from flat edge to flat edge, until you reach the bottom side. Write L for each rhombus you pass that leans left and R for each one that leans right. The picture shows a covering of a different board, where going down gives LRRLR.

\vspace{10pt}
\begin{center}
\begin{tabular}[b]{@{}c@{}}\fig{g45_Lrh}\\[3pt]\textbf{L}\end{tabular}\hspace{0.35in}%
\begin{tabular}[b]{@{}c@{}}\fig{g45_Rrh}\\[3pt]\textbf{R}\end{tabular}\hspace{0.9in}%
\fig{g45_ribbon_example}
\end{center}

\vspace{8pt}
Write the letters for each covering of the board from Problem 1. On the hexagon board from Problem 3, the top side has two short edges, so each covering gets two strings of letters. Call them its ribbons. Find exactly what one flip does to the ribbons.

\vspace{2pt}
\writelines{8}{\linewidth}

\newpage
\prob Explain why, whenever you make flips on the hexagon board and end at the covering you started with, you have made an even number of flips. Then explain why nobody can change covering E into covering F in Problem 3 with fewer flips than your answer.

\vspace{4pt}
\writelines{16}{\linewidth}

\newpage
\prob How many different coverings does the hexagon board from Problem 3 have? Draw or list them below, and explain how you know you have found them all.

\vspace{14pt}
\begin{center}
\foreach \r in {1,...,5}{%
\fig{g45_hex2_small}\hspace{0.25in}\fig{g45_hex2_small}\hspace{0.25in}\fig{g45_hex2_small}\hspace{0.25in}\fig{g45_hex2_small}\hspace{0.25in}\fig{g45_hex2_small}\par\vspace{14pt}}
\end{center}

\writelines{3}{\linewidth}

\newpage
\prob How many different ways can blue rhombuses \icon{B} cover this board? Explain how you know.

\vspace{14pt}
\begin{minipage}[t]{4.3in}
\vspace{0pt}
\fig{g45_b133}
\end{minipage}%
\begin{minipage}[t]{3.2in}
\vspace{0pt}
\writelines{12}{\linewidth}
\end{minipage}

\vspace{6pt}
\writelines{6}{\linewidth}

\end{document}
